\documentclass[11pt]{article}
\usepackage[a4paper,margin=24mm]{geometry}
\usepackage[T1]{fontenc}
\usepackage{lmodern,amsmath,amssymb,amsthm,mathtools,booktabs,longtable,array}
\usepackage[dvipsnames]{xcolor}
\usepackage{graphicx,microtype,fancyhdr}
\usepackage[colorlinks=true,linkcolor=MidnightBlue,urlcolor=MidnightBlue,citecolor=MidnightBlue]{hyperref}
\hypersetup{pdftitle={Passive Hamiltonian calibration from finite photon counts},pdfauthor={Astra research construction},pdfsubject={Constructive passive calibration, physical costs, and whole-CDF reliability certificates}}
\newtheorem{theorem}{Theorem}
\newtheorem{lemma}[theorem]{Lemma}
\newcommand{\Tr}{\operatorname{Tr}}
\newcommand{\supp}{\operatorname{supp}}
\newcommand{\id}{\mathrm{id}}
\newcommand{\ee}{\mathrm e}
\newcommand{\norm}[1]{\left\lVert #1\right\rVert}
\newcommand{\cL}{\mathcal L}
\newcommand{\cJ}{\mathcal J}
\newcommand{\cP}{\mathcal P}
\pagestyle{fancy}\fancyhf{}
\fancyhead[L]{\small Passive Hamiltonian calibration}
\fancyhead[R]{\small Astra scientific construction}
\fancyfoot[C]{\thepage}\setlength{\headheight}{14pt}
\title{\textbf{Passive Hamiltonian calibration\\from finite photon counts}}
\author{Astra research construction}\date{5 November 2026}
\begin{document}
\setlength{\emergencystretch}{3em}
\hypersetup{pageanchor=false}
\begin{titlepage}\maketitle\thispagestyle{empty}\vspace{7mm}
\noindent\textbf{Completed mathematical capability.}
Learn every coefficient of a static, bounded-degree local Hamiltonian from one continually evolving source, using only classical measurements of its calibrated outgoing Pauli fields. The source need not be prepared, controlled or read out. Calibration uses finite local oscillators and photon counters; it requires no coherent field collector or quantum memory. Observation windows stay fixed as coefficient accuracy improves.

\vspace{5mm}\noindent\textbf{Mechanism.}
Matched local oscillators give signed Poisson projection events. Time-reversing a finite-window correlation removes every even Hamiltonian order. A positive diagonal response, a global nonlinear inverse and a connected-set classical compiler complete the operation. A passive-output lower bound matches its precision and simultaneous-learning exposure dependence up to logarithms and fixed local/access factors.

\vspace{5mm}\noindent\textbf{Evaluated limits.}
Two million raw blocks were simulated for each of two finite-oscillator settings, including 40\% loss. The conservative uniform schedule at the earlier fine-precision target still requires approximately \(10^{26.17}\) blocks and \(10^{27.22}\) source-time units. These are sufficient bounds, not lower bounds, proposed device settings or an achieved practical advantage.

\vspace{5mm}\noindent\textbf{Independent completed companion.}
A finite certificate of the entire first-failure CDF is constructed for succinct reversible models with a priced static rare-event oracle, product equilibrium and equilibrium-safe initialization. Its confidence covers every positive rational deadline. A new changing-invariant-law inequality supplies acquired-model uncertainty transfer.
\vfill
\noindent The complete earlier acquisition laws, deferred first-input recovery, bounded-neighborhood construction, reliability sensitivity results and independent thermal-effusion operation are retained at their exact scopes. The new classical calibration does not remove the separate quantum archive's hardware or storage costs.
\end{titlepage}\hypersetup{pageanchor=true}
\tableofcontents\newpage

\section{Chosen question and the consequence actually earned}
The question is whether a many-body dynamical law can be learned through its emitted fields without first implementing a coherent reader for those fields or gaining direct access to the source. The completed answer is affirmative for a specified, experimentally substantive port model. The receiver records signed photon counts and reconstructs the coefficients of a bounded-degree Hamiltonian. It does not reconstruct an arbitrary quantum input from a classical record.

The distinction changes a broad access requirement. In the preceding deferred-calibration construction, known output purifiers supplied local probes, but obtaining those probes invoked a coherent continuous-output interface. Here a conventional destructive measurement supplies the needed classical response directly. Finite collection loss is allowed. Calibration can therefore be performed independently of the extremely costly P1 coherent archive. If that archive is also available, the preserved correlated-calibration theorem composes this learner with first-input recovery.

The new result is conditional on local Pauli fields, their calibrated phases, gains and rate, a complete interaction dictionary and stationary dynamics throughout the exposure. It is not general spectroscopy of arbitrary natural emissions. In particular, Hermitian Pauli jump operators are a substantive engineered or otherwise supplied coupling; an ordinary spontaneous-emission operator is not interchangeable with them.

The strongest relevant alternatives already contain substantial parts of the mathematics. Gu\c{t}\u{a}--Kiukas establish identifiability and local asymptotic statistics for ergodic quantum outputs \cite{guta}. Higher output moments and quantum polyspectra are established and have been fitted to experimental records \cite{hagele,sifft}. Finite-local-oscillator photodetection is also established \cite{banaszek,finiteLO}. None of these ingredients is claimed as a new discovery here.

Recent finite-response learning is an especially strong comparator. Arad et al. learn both Hamiltonian and dissipative coefficients with near-optimal experiment and evolution-time scaling, without a known interaction structure, using product-state source preparations and Pauli source measurements \cite{arad}. Song et al. give single-time projected response contraction as well as interpolation for local Lindbladians \cite{song}. Jim\'enez-Rodr\'iguez et al. use continuous weak records to learn time-dependent Hamiltonians with separable probe preparations \cite{jimenez}. The present result supplies a different acquisition interface: no such source preparations or source measurements, known noise rather than arbitrary unknown dissipators, and a complete bounded-degree candidate list. It does not dominate those methods where their stronger access is available.

The achieved additional connection is the full finite classical-output measurement, uniform inverse, local classical computation, finite detector error budget and access-matched passive lower bound. Historical priority for the complete connection has not been established. The evaluated costs do not establish a civilisation-changing practical capability.

\section{Source model and access contract}
Let \(S\) contain \(n\) qubits with a static Hamiltonian
\[
 H(\theta)=\sum_{P\in\cP}\theta_P P,\qquad
 1\le |\supp P|\le k,\qquad M=|\cP|.
\]
The distinct nonidentity Pauli dictionary \(\cP\) is known and complete. An identity coefficient is unobservable and is excluded. At most \(d\) candidates touch any site, and
\[
 \theta\in\Theta_J:=\left\{\theta:
 \max_i\sum_{P\ni i}|\theta_P|\le J\right\}.
\]
Fixed \(k,d\) are part of the polynomial classical-computation statement. Supplying all weight-\(k\) supports on an arbitrary graph may make \(d\) grow with \(n\); that is not silently a fixed-degree instance.

Each site has three independent calibrated vacuum-input ports
\[
 L_{i,a}=\frac{\sqrt r}{2}\sigma_{i,a},\qquad a=X,Y,Z,
\]
and the unconditional generator is
\[
 \cL_H(\rho)=-i[H,\rho]
 +r\sum_i\bigl(I_i/2\otimes\Tr_i\rho-\rho\bigr),\qquad r>0.
\]
The source begins in any state, possibly entangled with an inaccessible reference. Its maximally mixed state is invariant for every allowed \(H\). Loss from source to registered optical modes is calibrated, state-independent vacuum loss with efficiency \(\kappa>0\). Different known efficiencies can be normalized separately; a positive common floor bounds the resources.

The optical receiver supplies balanced interferometers, coherent phase references, photon counters, classical clocks and routing. Detection efficiency, port labels and axes, gains and phases are acquired separately or supplied with stated bounds. Constant photon flux does not certify the operator-valued port model, purity, independence, absence of hidden baths or completeness of the Hamiltonian dictionary. Unknown drift or omitted interactions are not removed by the inference theorem.

There are \(3n\) outgoing channels. A growing detector array is allowed and paid. The result does not grant a fixed aggregate optical passband, a microscopic completion of a white-noise field, or a calibrated detector response on frequencies that it cannot absorb. The actual spectral delivery and detection interface must meet the declared mode-loss and response-error contract.

\section{A finite optical measurement with bounded flux}
For a single port, interfere its detected field with a real coherent local oscillator of amplitude \(\beta\), then count both balanced output arms. The conditional jump operators are
\[
 J_\pm=\frac{\sqrt\kappa L\pm\beta I}{\sqrt2}.
\]
This operation takes place on the outgoing field; it leaves the unconditional source dynamics unchanged. Choose the matched setting
\[
 \beta=\frac{\sqrt{\kappa r}}2,
 \qquad J_\pm=\pm\sqrt{\frac{\kappa r}{2}}\,\Pi_\pm,
 \qquad\Pi_\pm=(I\pm\sigma)/2.
\]
Thus a registered click has a Pauli-projection conditional update and a sign. Its total rate is the scalar
\[
 \sum_\pm J_\pm^\dagger J_\pm=RI,\qquad R=\frac{\kappa r}{2}.
\]
The unmarked clocks of all ports are independent Poisson processes, for every source/reference state and every \(H\). Their marks carry the information. Lost photons remain part of the unconditional calibrated generator.

Choose adjacent windows \([0,w]\), \([w,2w]\), with \(q=rw\le1\). Let
\[
 m=Rw=\kappa q/2,
 \qquad D_{i,a}^{f,b}=N_{i,a,+}^{f,b}-N_{i,a,-}^{f,b},
 \qquad X_{i,a}^{f,b}=D_{i,a}^{f,b}/m.
\]
Every \(X\) is an ordinary classical count difference. The signed jump map is \(R\cJ_\sigma\), where
\(\cJ_A(X)=(AX+XA)/2\). Consequently products of distinct selected port-window variables have the same normalized regression moments as the ideal quadrature construction, exactly, at this finite oscillator power.

For target \(P\) of weight \(b\), select a site \(i\) and a Pauli \(B_i\) anticommuting with \(P_i\). The signed Pauli \(Q=iPB_i\) has the same support as \(P\). Form
\begin{align*}
 Y_P^\to&=\left(\prod_{j\in\supp P}X_{j,Q_j}^{f}\right)X_{i,B_i}^{b},\\
 Y_P^\leftarrow&=X_{i,B_i}^{f}\left(\prod_{j\in\supp P}X_{j,Q_j}^{b}\right),
 &Y_P&=(Y_P^\to-Y_P^\leftarrow)/2.
\end{align*}
A known overall Pauli sign is applied classically. In each product, all factors use distinct ports: even the repeated site uses different axes \(Q_i\ne B_i\). All \(M\) products are obtained from the same \(6n\) window counts. Their evaluation costs \(O(kM)\) arithmetic per block; there are not \(M\) separately acquired source batches.

For \(s=b+1\) factors, let \(T_j\) be their total unsigned counts. They are independent \(\operatorname{Pois}(m)\), and \(|D_j|\le T_j\). Hence, uniformly in every conditional source state,
\[
 \mathbb E(Y_P^\to)^2\le
 \prod_{j=1}^s\frac{\mathbb E T_j^2}{m^2}
 =\left(1+\frac2{\kappa q}\right)^s=:V_b.
\]
The same variance bound holds for the odd difference, by
\((A-B)^2/4\le(A^2+B^2)/2\).

The more general finite-oscillator observation is also valid. If \(\mu\) is the mean detected oscillator photon number per window, a balanced count difference divided by \(\sqrt{\kappa\mu q}\) gives the same normalized first moment and all distinct-window cross moments. A useful general second-moment bound is
\[
 V_b^{\rm general}=
 \left(\frac2{\kappa q}+1+\frac1{4\mu}\right)^{b+1}.
\]
For broadband counting the exact mode identity contains the \emph{total} signal photon number: \(\mathbb E_{\rm LO}D^2/\mu=X_{\rm matched}^2+N_{\rm total}/\mu\). Replacing it by only the matched-mode number would be incorrect. The preferred matched setting \(\mu=\kappa q/4\) uses the sharper independent-total-count bound above. It guarantees bounded flux, not the smallest observed variance on every device.

\section{Time reversal makes the response invertible}
Write \(G_P(H)=\mathbb E_{I/2^n}Y_P\). Ordered output regression inserts Jordan maps between source propagators. In the normalized Hilbert--Schmidt inner product,
\[
 \cL_H^\dagger=\cL_{-H},\qquad \cJ_A^\dagger=\cJ_A.
\]
Reverse the trace word and time coordinates \(t\mapsto2w-t\). This exchanges the forward and reverse patterns and proves
\[
 \mathbb E Y_P^\leftarrow(H)=\mathbb E Y_P^\to(-H).
\]
All even orders in the Hamiltonian therefore cancel exactly. This is a stationary identity; the nonstationary acquisition bias is bounded separately.

At first order, a nonzero final Pauli trace requires the inserted interaction to be precisely \(P\). All such contributions have the same sign. With times measured in units of \(w\),
\begin{equation}
 G_P(H)=wC_b(q)\theta_P+\mathcal R_P(H),
 \label{eq:response}
\end{equation}
where the remainder begins at cubic order and
\begin{align}
 C_b(q)&=2\int_1^2\!dz\int_0^1\!dx\,
 \ee^{-q(z-x)}\int_x^z\!du\,[A_q(u)]^{b-1},\label{eq:kernel}\\
 A_q(u)&=\int_0^1\ee^{-q|u-t|}\,dt.\nonumber
\end{align}
The insertion lies between the two measurements at site \(i\). Its noise length there is \(z-x\); at each other site it is \(|u-t|\). Integrating those front times gives \eqref{eq:kernel}. Positivity yields
\[
 2=C_b(0)\ge C_b(q)\ge c_b:=C_b(1)>0,
 \qquad C_b(q)\ge2\ee^{-2bq}.
\]
For small \(q\), the entirely rational lower bound \(2(1-2bq)\) is also useful. Numerical adjacent-window values are
\[
 c_1=0.6681289718,\quad c_2=0.4041133844,\quad
 c_3=0.2592800494,\quad c_4=0.1736661755.
\]
They are quadrature diagnostics of the explicit integral; interval quadrature or the analytic lower bounds supply finite certified normalizations.

The order \(b+1\) has a precise necessity interpretation. In the family \(H_\theta=\theta P\), every stationary local quadrature moment of order at most \(b\) is invariant under \(\theta\mapsto-\theta\). An odd Hamiltonian word with nonzero trace needs the measured single-site Pauli product to equal \(P\). Fewer than \(b\) factors cannot cover its support. Exactly \(b\) factors force one matching Pauli per site, so every intermediate product commutes with \(P\) and the first commutator vanishes. This is a statement about moment order, not a lower bound against arbitrary processing of a complete record.

\subsection{A uniform inverse at every positive emission rate}
A weight-\(v\) Pauli has noise damping \(rv\), while its outgoing absolute commutator weight is at most \(2Jv\). Put
\[
 a=\max(1,k-1),\quad x=4aJw,\quad p_b=b/a.
\]
After \(\ell\) nonzero Hamiltonian insertions, support has size at most \(b+(k-1)\ell\). The combined time simplex has length at most \(2w\). Discarding damping gives the absolute order bound
\[
 |G_{P,\ell}(H)|\le\frac{(p_b)_\ell}{\ell!}x^\ell.
\]
For two parameters in \(\Theta_J\), telescope one insertion and use
\(\Delta J\le d\norm{\theta-\theta'}_\infty\). Summing odd orders \(\ell\ge3\) gives a normalized remainder Lipschitz bound
\begin{equation}
 \rho_{\rm time}=\max_{b\le k}\frac{4db}{C_b(q)}
 \left[\frac{(1-x)^{-p_b-1}+(1+x)^{-p_b-1}}2-1\right].
 \label{eq:rho}
\end{equation}
Any \(w\) with \(q\le1,x<1,\rho_{\rm time}\le1/4\) suffices. In particular, set
\[
 C_*:=\max_{b\le k}b/c_b,\quad
 A_k=64a\sqrt{C_*},\quad
 w=(r+A_k\sqrt d\,J)^{-1}.
\]
Then \(x\le1/(16\sqrt{dC_*})\), \(p_b\le2\), and the bracket in \eqref{eq:rho} is at most \(16x^2\), proving the required bound. The window depends on dynamical and local-degree bounds, not on target accuracy or source size at fixed local parameters. Weak emission still raises \(V_b\) through \(q\).

When \(w=1/r\), retaining damping gives a sharper certificate. Extend each waiting time before a Hamiltonian insertion to infinity. Its integrated absolute transition matrix has column sum at most \(\alpha=2J/r\). The \(b\) regression gaps give order bound \({\ell+b-1\choose b-1}\alpha^\ell\), and
\[
 \rho_{\rm damp}=\max_{b\le k}\frac{2db}{c_b}
 \left[\frac{(1-\alpha)^{-b-1}+(1+\alpha)^{-b-1}}2-1\right].
\]
It is quadratic in \(J/r\). For \(k=2,d=12\), a scalar numerical evaluation places the sufficient \(\rho=1/4\) threshold at \(r/J\simeq106.83\). This is a sufficient contraction threshold, not an identifiability phase transition.

\subsection{Finite nonlinear reconstruction}
Let \(\lambda_P=wC_{|P|}(q)\), \(\lambda_*\) a positive certified lower bound on their minimum, and \(D_\lambda\) the diagonal matrix. For target coefficient error \(\eta\), set
\[
 \delta=\lambda_*\eta/32.
\]
Acquire \(y\) with \(\norm{y-G(\theta_*)}_\infty\le\delta\). A finite classical evaluator supplies \(\widetilde G\) with error at most \(\delta\). Starting at zero, iterate
\[
 z_t=\theta_t+D_\lambda^{-1}\bigl(y-\widetilde G(\theta_t)\bigr),
 \qquad\theta_{t+1}\in\arg\min_{\theta\in\Theta_J}\norm{\theta-z_t}_\infty.
\]
The nearest feasible projection is a linear program using auxiliary absolute-value variables. It need not be nonexpansive: feasibility of \(\theta_*\) gives the valid factor two. Thus
\[
 e_{t+1}\le\tfrac12e_t+4\delta/\lambda_*,
 \qquad e_t=\norm{\theta_t-\theta_*}_\infty.
\]
After \(\lceil\log_2(2J/\eta)\rceil_+\) steps, the ideal arithmetic error is at most \(3\eta/4\). The remaining allowance covers finite diagonal normalization, approximate LP projection and arithmetic. A relative diagonal error \(O(\eta/J)\) suffices; all these quantities have polynomial bit length. The cubic remainder is corrected through the nonlinear response, not declared negligible.

\subsection{A classical forward compiler, without exponential source simulation}
Literal enumeration of \(\ell\) nested-commutator words can cost \((d\ell)^\ell\), which is only quasipolynomial when \(\ell=O(\log(1/\delta))\). A connected-set construction avoids this gap.

Root the candidate-interaction overlap graph at the measured support of \(P\). For a finite interaction set \(S\), define
\[
 \Phi_P(S)=\sum_{T\subseteq S}(-1)^{|S|-|T|}G_P(H_T),
 \qquad H_T=\sum_{Q\in T}\theta_QQ.
\]
Disconnected components cancel. M\"obius inversion partitions every Dyson word by its distinct interaction set: it appears once in \(\Phi_P(S)\) exactly when that set is \(S\). Truncating root-connected sets at size \(L\) therefore omits only orders exceeding \(L\). Either absolute tail
\[
 \sum_{\substack{\ell>L\\\ell\ \mathrm{odd}}}
 \frac{(p_b)_\ell x^\ell}{\ell!},\qquad
 \sum_{\substack{\ell>L\\\ell\ \mathrm{odd}}}
 {\ell+b-1\choose b-1}\alpha^\ell
\]
controls the error in its certified regime.

There are at most \((4kd)^L\) root-connected sets of size \(L\), via plane rooted spanning trees and neighbor labels. Each occupies at most \(b+(k-1)L\) qubits. For each subset response, a register \(M_U\) for each subset of measured channels satisfies, during a dimensionless window,
\[
 \frac{dM_U}{ds}=w\cL_{H_T}M_U+
 \sum_{j\in U}\cJ_{Q_j}M_{U\setminus\{j\}}.
\]
The initial empty register is maximally mixed; other registers vanish. Propagate the front and back windows and take the trace, then reverse their roles. A finite matrix exponential action computes this system. At most \(2^{b+1}4^{b+(k-1)L}\) operator coordinates are needed.

Choose \(L=O(\log(1/\delta))\) to bound the tail by \(\delta/2\). If \(N_L\) sets are retained, evaluate each subset response to \(\delta/(2N_L2^L)\). This pays the cancellation precision. Matrix norms grow polynomially with \(L\) at fixed parameters, so certified exponential actions have finite polynomial bit cost. Overall classical time and memory are polynomial in \(n,1/\eta\) at fixed \(k,d\) and access ratios, although their exponents and constants may be large.

The combinatorial compiler was implemented with rooted-set enumeration, cached subset responses, local-support compression and sparse Pauli-basis moment propagation. One three-qubit, five-term diagnostic at \(L=3\) used 26 and 20 rooted sets for two targets. Its responses differed from an independent full density-matrix calculation by \(7.33\times10^{-8}\) and \(1.37\times10^{-7}\). The conservative analytic tails were much larger: 0.0548 and 0.00863. Maximum augmented dimension was 256; the saved evidence records 196 sparse exponential actions and actual matrix/right-hand-side application counts. These are floating-point \texttt{expm}/\texttt{expm\_multiply} calculations, not the theorem's certified finite-arithmetic backend. The two-qubit sampled-data diagnostic below uses its cheaper full small-system response solver.

\section{Acquisition, finite service and the complete cost bill}
\subsection{Conditional concentration from one continuing source}
The normalized regression effect has operator norm at most one. Each Jordan insertion and each adjoint source channel is a contraction, and the normalized integration volumes sum to one. A sufficient entropy-forgetting bound is
\[
 \norm{\rho(t)-I/2^n}_1
 \le2\sqrt{n\log2}\,\ee^{-rt/2}.
\]
The Hamiltonian preserves entropy relative to the maximally mixed state; local depolarization supplies the decay. Conditional on all previous classical results, the next source state is still a density matrix, so the same bound applies. Source entanglement with an old quantum archive does not invalidate this marginal conditional-mean argument.

Before every two-window experiment discard output for
\[
 T_{\rm burn}=\frac2r\log\frac{16\sqrt{n\log2}}{\delta},
\]
or zero if the displayed value is negative. This makes the mean bias at most \(\delta/8\). The source is not assumed to have mixed because a short record looks stationary. Nor is an \(O(1/N)\) approximation to the entire record distribution needed: conditional concentration bounds the mean directly.

Use a finite counter cap \(K\). A coefficient's observation is zero if any of its at most \(2(k+1)\) selected modes overflows. Unrelated overflows do not discard the block. With
\[
 V=\max_{b\le k}(1+2/(\kappa q))^{b+1},\qquad
 B=\max_{b\le k}(K/m)^{b+1},
\]
the clipping bias is at most
\(\sqrt{2V(k+1)\Pr(\operatorname{Pois}(m)>K)}\).
Choose the tail below \(\delta^2/[128V(k+1)]\). This is a finite scalar computation; the rational bound
\((3m/(K+1))^{K+1}\) is sufficient when \(K+1\ge m\).

Allocate at most \(\delta/8\) each to forgetting, clipping and physical moment bias, and \(\delta/2\) to statistical fluctuation. A conservative sufficient count is
\begin{equation}
 N=\left\lceil128\left(\frac V{\delta^2}+\frac B\delta\right)
 \log\frac{2M}{\zeta}\right\rceil.
 \label{eq:shots}
\end{equation}
Conditional Bernstein--Freedman concentration and a union bound imply simultaneous response error at most \(\delta\) with probability \(1-\zeta\). Empty windows and every failed/overflow outcome count toward \(N\). Independence between different coefficient estimators is not assumed. Small bounded gain/timing distortions can use conservatively enlarged variance constants; the displayed allocation has slack for a factor-two enlargement.

The window duration is independent of \(\eta\). Timing precision, gain/phase precision and the number of blocks are not. At fixed \(k,d,J/r,\kappa\), the leading precision dependence is \(\eta^{-2}\) up to logarithms. In a certified \(w=1/r\) regime,
\[
 N=\widetilde O_{k,d,\kappa}\!\left(\frac{r^2}{\eta^2}
 \log\frac M\zeta\right),\qquad
 T_{\rm acq}=\widetilde O_{k,d,\kappa}\!\left(\frac r{\eta^2}
 \log\frac M\zeta\right).
\]
The tilde includes paid forgetting and finite-counter factors. This statistical source-exposure law counts the waiting and observation intervals; any additional accuracy-dependent receiver service or reference calibration must be added as below. It is not an optimality claim for the complete hardware bill. The weak-port factor \((rw)^{-(k+1)}\) is explicit in \(V\), and is not claimed optimal.

\subsection{Detector dead time, timing and clocks}
Assume downstream dead time only deletes events, with maximum recovery time \(d_0\). Split each arm equally into \(L\) labelled detector rails, then add signed counts. A guard of at least \(d_0\) handles arbitrary initial detector blindness. Same-rail close-arrival pairs dominate missed photons. For one selected window their expected number is at most
\[
 \mathbb EC_j\le R^2(w+d_0)d_0/L.
\]
The \(d_0\) prefix pays arrivals immediately before the window. Since every factor in one product uses a distinct port, \(C_j\) is independent of the other total counts. Telescoping the \emph{unclipped} products gives
\begin{align*}
 \mathbb E\left|\prod_jD_j/m-\prod_j\widetilde D_j/m\right|
 &\le\sum_j\mathbb E(C_j/m)\prod_{\ell\ne j}\mathbb E(T_\ell/m)\\
 &\le (b+1)\frac{Rd_0}{L}\left(1+\frac{d_0}{w}\right).
\end{align*}
The odd average obeys the same bound. Clip the physical observations afterward: deletion only reduces their total counts, so the same overflow and variance domination still apply. This ordering avoids multiplying the service error by the large counter cap \(B\).

There are \(6nL\) physical counters in the literal implementation. Rail routing loss belongs in \(\kappa\); rail lengths, signal propagation, synchronization, recovery, readout and aggregation belong in the actual execution cost. Afterpulses, arbitrary dark counts or correlated backgrounds are different detector models and need their own calibrated bias/variance allocation.

For deterministic bounded window-endpoint errors, or independent calibrated timing jitter, the normalized signed-moment kernel has norm at most one. If each endpoint moves by at most \(\Delta t\), its normalized integration kernel changes in \(L^1\) by at most \(2\Delta t/w\). Tensor telescoping bounds the product bias by \(2s(\Delta t/w)(1+2\Delta t/w)^{s-1}\), where \(s=b+1\); for \(\Delta t\le w/(8s)\) this is at most \(4s\Delta t/w\). Apply this before clipping. A conservative allocation is \(\Delta t\le\delta w/[128(k+1)]\). Pad the discard interval so its full duration precedes the earliest possible accepted event. Outcome-dependent timing errors do not follow from this argument. Actual widths enlarge total-count variance by at most \((1+2\Delta t/w)^{2s}\).

With both endpoint uncertainty and dead time, a conservative combined deletion-bias bound is
\[
 \frac{sRd_0}{L}\left(1+\frac{2\Delta t+d_0}{w}\right)
 \left(1+\frac{2\Delta t}{w}\right)^{s-1}.
\]
Choose the integer rail count using this expression when both imperfections are present. It reduces to the previous dead-time expression at zero timestamp error; the correction is negligible at the illustrative tolerance below.

The array construction resolves detector busy time; it does not supply arbitrarily precise timestamps from a detector with fixed unknown jitter. The stated timing bound must be physically available and calibrated. A direct digital timestamp implementation uses clock period at most the allocated tolerance, hence clock frequency of order \(1/\Delta t\), plus enough bits for the complete exposure. Actual detector latency uncertainty and bandwidth must meet the same timing contract. A fixed nonvanishing jitter floor requires a different, explicitly modeled response inversion; such a general filter decoder is not claimed here. The protocol therefore retains growing timing/control demands and makes no fixed-full-hardware completion claim.

The source-normalized insertion is a bounded map. Uniform local gain, quadrature-phase or response-map errors therefore accumulate at most linearly in the \(b+1\) factors to first order, with an explicit finite product bound. Port calibration must meet the allocated moment-error tolerance. Rate \(r\) can be estimated from unmixed counts at rate \(\kappa r/4\) if \(\kappa\) is independently calibrated, paying Poisson exposure and systematic errors. Those counts do not identify \(r\) and \(\kappa\) separately if both are unknown, and do not certify the full port operator model.

\subsection{Exposure, illumination, memory and computation}
The source must remain within its static-model bounds throughout
\[
 T_{\rm acq}=N\bigl(T_{\rm burn}+2w+T_{\rm service}\bigr),
\]
where service includes any additional nonoverlapped guard, ramp, acquisition and routing latency. Known fixed propagation can be compensated by classical timestamps, with its precision paid. The oscillator's total detected flux while on is \(3n\kappa r/4\), independent of \(w\). Its detected photons in the counted windows total
\[
 \frac{3n\kappa rwN}{2}.
\]
Leaving it on throughout forgetting instead costs \(3n\kappa rT_{\rm acq}/4\). Incident photons divide these values by the actual oscillator path/detection efficiency \(\eta_{\rm LO}\). Gating can save illumination only after paying its finite slew, calibration and ramp photons. Neither vacuum preparation nor the source-port engineering is supplied by these optical photon counts.

Records can be streamed. A block needs \(O(n\log K)\) count storage and \(O(kM)\) product work; running sums and flags use finite \(O(M)\) registers whose bit length grows with \(\log N\) and target precision. No quantum memory is needed for calibration. The nonlinear inverse adds its connected-set computations, finite matrix actions, numerical precision and projection LPs. The literal forward compiler can still be expensive, even though polynomial at fixed local degree.

\section{A passive-exposure lower bound}
The standard inverse-square precision rate is not merely an artifact of photon counting in this access model. Consider one qubit with \(H_\theta=\theta Z\) and known depolarization rate \(r\), initially maximally mixed. Strengthen the receiver by giving it all output fields, the final source and an initial Bell purifier. Its global final state is pure. For parameter separation \(\Delta\), the pure-state overlap satisfies the exact two-dimensional cross-evolution equation
\[
 a''+ra'+\Delta^2a=0,\qquad a(0)=1,\quad a'(0)=0.
\]
For \(|\Delta|\le r/2\), with \(s=\sqrt{r^2-4\Delta^2}\),
\[
 a(T)=\ee^{-rT/2}\left[\cosh(sT/2)+(r/s)\sinh(sT/2)\right]
 \ge\ee^{-2\Delta^2T/r}.
\]
The limit at \(s=0\) is continuous. Helstrom discrimination therefore has error
\[
 p_e=\frac{1-\sqrt{1-a(T)^2}}2\ge\frac14\ee^{-4\Delta^2T/r}.
\]
Take \(n\) independent coefficients \(\pm2\eta\), so \(\Delta=4\eta\). Their strengthened output ensemble factorizes. Optimal whole-string guessing probability is exactly \((1-p_e)^n\): product Helstrom measurements achieve it, and tensoring the one-qubit discrimination dual certificates gives the upper bound. Any simultaneous \(\eta\)-accurate learner with success at least \(1-\zeta\) must therefore obey
\[
 T\ge\frac r{64\eta^2}
 \log\frac{n}{4[-\log(1-\zeta)]},\qquad\eta\le\min(J/2,r/8),
\]
whenever the right side is positive. This gives \(\Omega((r/\eta^2)\log(n/\zeta))\) in its asymptotic regime. The exact one-qubit pure-state Fisher information is
\[
 8\left[\frac Tr-\frac{1-\ee^{-rT}}{r^2}\right]\le8T/r.
\]
The comparison permits receivers stronger than the construction but fixes passive source dynamics. Source interventions, changing the bath couplings, active error correction or a different parameter family can change the lower bound. It establishes the precision and simultaneous-learning exposure comparison, not optimal loss dependence, weak-port constants or classical runtime. Stationary-output local asymptotic normality is prior machinery \cite{guta}; this is a concrete finite hard-family calculation for the stated access.

\section{Executed evidence and evaluated sufficient costs}
\subsection{Actual photon-count trajectories}
The saved program simulates the event-driven quantum trajectory directly. Scalar total jump intensity permits exact Poisson event counts and sorted uniform event times. Between events it applies the exact two-qubit Hamiltonian unitary. Detected signs are sampled from their Born probabilities; unobserved loss jumps remain in the source dynamics. No time-step approximation is used.

The Hamiltonian is
\[
 H=0.37XI-0.23IY+0.41ZZ+0.19XY-0.17ZX.
\]
All 15 nonidentity two-qubit Paulis are estimated, including absent terms. Both runs use \(r=4,w=0.25,\kappa=0.6\), eight inverse iterations, and two million raw blocks. The stronger-oscillator run uses \(\mu=1\); the matched run uses \(\mu=0.15\). The random seed and full numerical records are preserved.

\begin{center}\small
\begin{tabular}{lrr}\toprule
Recorded quantity & Stronger LO & Matched LO\\\midrule
Raw two-window blocks & 2,000,000 & 2,000,000\\
Simulated source + LO events & 29,992,217 & 9,596,361\\
Maximum coefficient error & 0.03070 & 0.05944\\
Largest absolute response residual / SE & 1.914 & 3.021\\
Source observation exposure & \(10^6\) & \(10^6\)\\\bottomrule
\end{tabular}\end{center}
The exposure entries exclude burn-in and service. Computational runtime is not a physical exposure claim. Starting each simulated block at \(I/4\) supplies the stationary input numerically; the physical waiting protocol is not simulated. The simulation includes finite oscillator power and loss, but not dead time, clock errors, dark counts, finite clipping, factories or a device's spectral response.

The exact regression calculation satisfies the reversal identity to \(4.17\times10^{-17}\). With exact response means, the inverse error falls from \(9.43\times10^{-3}\) to \(3.06\times10^{-11}\) in eight steps. With sampled means it reaches a statistical floor. This small diagnostic applies the response-correction iteration without the theorem's feasible-set projection. These are floating-point diagnostics, not formal enclosures or an execution of the whole uniform schedule. The useful \(w=0.25\) point is outside the conservative uniform contraction certificate; its convergence does not certify the whole promised parameter domain at that window.

\begin{figure}[htbp]\centering
\includegraphics[width=\textwidth]{next2_root/figures/photon_count_calibration.png}
\caption{Executed synthetic photon records and nonlinear inversion. Both oscillator settings, including the larger matched-setting error, are retained. The exact-mean curve isolates the deterministic inverse from sampling noise. This small example does not execute the generic connected-set compiler or the worst-case physical schedule.}
\end{figure}

\subsection{A uniform certificate at the previous fine-precision target}
For an explicit comparison with the preserved deferred-recovery schedule, use
\[
 n=2,\ M=15,\ k=2,\ d=12,\ J=1.14,\ r=4,\ \kappa=0.6,
 \qquad \eta=1.4526409921\times10^{-5}.
\]
Choose \(w=3/500=0.006\), \(q=0.024\), and the matched oscillator \(\mu=0.0036\). The rational bound \(C_b(q)\ge2(1-2bq)\) gives
\[
 \lambda_*\ge0.010848,\qquad
 \rho_1\le0.056686,\quad\rho_2\le0.238930<1/4.
\]
These contraction comparisons were performed with exact rational arithmetic. With \(\zeta=0.001\), the resulting allocations are
\begin{center}\small
\begin{tabular}{lr}\toprule
Sufficient resource or tolerance & Evaluated value\\\midrule
Normalized response accuracy \(\delta\) & \(4.92445\times10^{-9}\)\\
Uniform second-moment bound \(V\) & \(2.73747\times10^6\)\\
Finite count cap \(K\) & 9\\
Raw shared blocks & \(10^{26.1731}\)\\
Burn duration per block & 11.0325\\
Source exposure before extra service & \(10^{27.2162}\)\\
Illustrative endpoint timing tolerance & \(7.69446\times10^{-14}\)\\
Inverse iterations, conservatively rounded & 19\\\bottomrule
\end{tabular}\end{center}
Times are model units; no device clock has been assigned to them. Logarithmic resource evaluations used 90-digit arithmetic and conservative rounding conventions, rather than a machine-checked interval library.

For illustration only, fixing detector dead time at \(d_0=10^{-6}\) model units and allocating \(\delta/32\) to it gives about \(2.34\times10^4\) rails per port from the proved deletion-only bound. Different actual dead time, routing loss, propagation and spectral response change this bill. The illustrative timing and phase/gain demands remain severe. At oscillator efficiency equal to \(\kappa\), counted-window incident LO allocation is about \(10^{25.03}\) photons; leaving it on through burn-in raises that to about \(10^{27.99}\), before extra service.

The preceding source/export measurement schedule required about \(10^{39.12}\) total raw trials at this coefficient target, before its coherent receiver compilation. This new sufficient schedule reduces that evaluated count by about thirteen orders of magnitude and changes the measurement hardware. The old two-site construction also had a no-burn-in advantage, which remains preserved. The protocols have different observations and access costs; this comparison establishes neither an optimality gap nor competitive laboratory performance. Both sufficient counts are far from lower bounds. The two-million-block simulations use a much less stringent achieved accuracy and must not be substituted for this uniform guarantee.

\section{Composition with the preserved first-input recovery theorem}
Suppose the accepted P1 coherent archive captures the first arbitrary input for time \(T_0\), at a supplied sufficient margin \(r>\beta_0\). Apply the new classical learner only to subsequent output. The Hamiltonian stays the same. From coefficient accuracy \(\eta\),
\[
 \norm{H-\widehat H}\le M\eta.
\]
The preserved correlated-calibration composition permits the classical calibration record to remain correlated with the archived quantum state. Its complete channel error includes
\[
 2\sqrt{n\ee^{-(r-\beta_0)T_0}}
 +2T_0M\eta+2\zeta
 +\epsilon_{\rm archive}+\epsilon_{\rm compilation}.
\]
The statistical success statement is uniform over the preceding source/archive input; no unjustified replacement by an independent calibration sample is made. The known-Hamiltonian reverse-reset decoder and its comparison theorem remain the accepted operation, not a newly rederived decoder here.

The archive must survive for the full acquisition time, classical inference time, routing delay and decoder execution. Its lossless coherent delivery, protected memory, finite compilation and growing hardware remain separate costs. The latest bounded-component P1 result pays a growing array and precision while retaining its fixed-response, spectral and access-order limitations. Its \(>10^{112}\)-cell sufficient one-qubit recipe is not an inexpensive buffer. Allowing loss in this later classical calibration does not weaken the archive's same-source first-ingress loss converse.

\section{Independent companion: a finite certificate for the whole failure CDF}
The independently developed reliability route completed a different consequential endpoint: one finite object containing relative first-failure probabilities at every positive rational deadline. The earlier package's computational proposal was explicitly incomplete. The following theorem is its new closure; it does not retroactively change that earlier status.

\begin{theorem}[Finite simultaneous risk curve]
Let a known undirected all-terminal network have \(m\) binary edges, rational product invariant law \(\pi\), failure set \(B\) and \(0<p=\pi(B)\le1/2\). Supply a succinct reversible continuous-time generator, ordinary local-update and failure-test algorithms, an exit-rate upper bound and a certified gap \(\gamma>0\). Start from \(\pi(\cdot\mid B^c)\). Supply the existing all-terminal static FPRAS at its actual bit cost.

For \(\epsilon,\delta_{\rm conf}>0\), finite preprocessing constructs a certificate which, with probability at least \(1-\delta_{\rm conf}\), returns a containing interval of ratio at most \(1+\epsilon\) for the first-failure CDF at every positive rational deadline. Work is polynomial in the succinct input length, the numerical inverse uniformized gap, \(1/\epsilon\) and \(\log(1/\delta_{\rm conf})\), including static-oracle work. It has no inverse global-failure-probability or mean-lifetime factor. Queries use polynomial bit work in certificate and deadline length, without new trajectories.
\end{theorem}

Product equilibrium does not require independent dynamics: reciprocal bounded-block transitions may accompany an operating independent refresh backbone. The gap floor and static access are supplied with their costs. A general failure predicate needs its own suitable static oracle; no FPRAS for arbitrary predicates follows. The initial equilibrium-safe law is a premise with a separate physical preparation cost.

\subsection{Spectral representation and the improved tail}
Choose a rational uniformization rate \(R_u\ge2\max_x[-Q(x,x)]\). Then \(P=I+Q/R_u\) is reversible and positive semidefinite. Write \(g=\gamma/R_u\), \(u=R_ut\), \(A=B^c\), \(a=1-p\), and \(K=I-P_A\). Genuine Poisson uniformization preserves the actual continuous-time path law, including null events.

Let \(j=\sum_{x\in A,y\in B}\pi(x)P(x,y)\). The safe-start CDF is
\[
 F(u)=\sum_i\alpha_i(1-\ee^{-\lambda_i u}),
 \qquad\alpha_i\ge0,\quad\sum_i\alpha_i=1.
\]
Poincar\'e, Rayleigh quotients and principal-submatrix interlacing give
\[
 gpa\le j\le p,\quad gp\le\lambda_1\le j/a\le p/a,
 \quad\lambda_2\ge g.
\]
Let \(r_f=1-\alpha_1\), \(\lambda=\lambda_1\). A normalized killed ground eigenvector extended by zero has full-chain variance \(p+ar_f\). Therefore
\[
 \lambda\ge g(p+ar_f),\qquad r_f/\lambda\le1/(ag).
\]
It follows that, with \(f(u,\lambda)=1-\ee^{-\lambda u}\),
\begin{equation}
 f(u,\lambda)\le F(u),\qquad
 \frac{F(u)-f(u,\lambda)}{f(u,\lambda)}
 \le\frac{r_f}{\ee^{\lambda u}-1}\le\frac2{gu}.
 \label{eq:tail}
\end{equation}
This is relative accuracy of failure probability, not relative accuracy of the tiny survival tail. The positive mixture also makes \(F\) increasing and concave, with \(F(t)/F(s)\le t/s\) for \(0<s\le t\).

\subsection{Bounded observations replace lifetime trajectories}
Draw \(y\sim\pi(\cdot\mid B)\), perform one \(P\) update, and accept if it enters \(A\). Acceptance is \(c=j/p\ge g/2\). Reversibility makes the accepted reversed edge \(x\to y\) a stationary boundary crossing.

Draw \(M\sim\operatorname{Pois}(u)\), place that crossing uniformly among \(M+1\) positions, and simulate the remaining \(M\) steps backward and forward with ordinary \(P\). If \(N_\times\) counts all \(A\to B\) crossings in the resulting path, return
\[
 W=\frac{1_{\{X_0\in A\}}}{N_\times}\in[0,1].
\]
A skeleton with \(\ell\) potential events has proposal density \(N_\times/(ju)\) relative to its unforced stationary probability: the factors \(\ell\) from the Poisson law and crossing placement cancel. Thus
\[
 aF(u)=pcu\,\mathbb EW,\qquad
 \mathbb EW\ge\frac{ag}{1+gpu}.
\]
This is an explicit boundary-Palm inverse-multiplicity operation, using established rare-event ideas \cite{karpluby}. Given a narrow static interval \([p_L,p_U]\), use it when \(ugp_U\le1\), so the mean is at least \(g/4\). Otherwise ordinary safe-start simulation has failure probability above \(1/3\). Every fitted middle-range observable then has an inverse-polynomial mean.

For the slow mode, start at \(\pi\mid B\), take one update, and return \(Z=1\) if it enters safety and stays safe for a further Poisson number of steps with mean \(\ell\). Then
\[
 p\mathbb EZ=a\sum_i\alpha_i\lambda_i\ee^{-\ell\lambda_i}.
\]
Put \(\xi=\epsilon/40\),
\[
 \ell=\left\lceil\frac{\lceil\log_2(4/(g\xi))\rceil}{g}\right\rceil,
 \qquad p_{\rm gate}=\xi/(2\ell).
\]
If \(p\le p_{\rm gate}\), spectral separation gives
\[
 (1-2\xi)\lambda\le p\mathbb EZ/a\le(1+\xi)\lambda,
 \qquad\mathbb EZ\ge g/4.
\]
The reverse window is set by the gap and accuracy, not by the lifetime.

\subsection{One finite object covers all deadlines}
There are three parts. For \(0<u\le\xi\), use the relative linear bound
\[
 (1-u/2)(j/a)u\le F(u)\le(j/a)u.
\]
There is no infinite grid accumulating at zero. In the middle, estimate the bounded observables on a finite geometric grid, then interpolate by monotonicity and concavity. In the rare branch the upper grid horizon is
\[
 H=\max\{\lceil2/(g\xi)\rceil,\ell+1\}.
\]
Beyond it, \eqref{eq:tail} and the sampled interval for \(\lambda\) give the entire tail. In the nonrare branch, \(p_L>p_{\rm gate}/2\); a polynomial horizon \(H=O(g^{-2}\xi^{-1}\log(1/(g\xi))\log(1/\xi))\) makes \(F\ge1-\xi\), so the remaining interval is \([1-\xi,1]\).

All samplers are finite. Sequential edge conditioning uses existing static all-terminal counting on contracted/deleted prefix graphs. For one conditional failure draw it makes at most \(2m\) static calls at explicitly requested accuracy. A positive mass has a bit-length floor \(\prod_i\min(q_i,1-q_i)\); its small value increases arithmetic precision, not a reciprocal-probability loop. The generic static FPRAS is imported machinery \cite{karger,cen}, not newly implemented here.

The Poisson law is truncated at
\[
 K_u=\lceil2u\rceil+\lceil\log_2(1/\tau_P)\rceil.
\]
Its omitted tail is at most \(\tau_P\). For rational \(u=A/B\), conditional probabilities on \(0,\ldots,K_u\) are sampled exactly using integer weights
\[
 A^j B^{K_u-j}K_u!/j!.
\]
Their construction, memory and random bits are charged. Approximate conditional-failure sampling is allocated as a bounded-mean bias. Boundary rejection, safe preparation and integer rejection have finite caps with separately allocated abort probabilities. The full proof specifies these allocations and the simultaneous confidence event.

For arbitrary binary rational deadlines, tiny \(1-\ee^{-x}\) values are evaluated by outward rational Taylor bounds after a power-of-two scaling, then integer powering. Large arguments are saturated with an explicit outward interval. This avoids subtracting two approximate floating-point numbers near one. Grid size and bit precision remain polynomial. No time-length simulation is hidden in an astronomical deadline query.

The resulting worst endpoint ratio is at most \(((1+\xi)/(1-\xi))^9\le1+\epsilon\). The discrete companion uses actual positive-semidefinite endpoint kernels, exact step counts and outward repeated squaring of \(1-(1-\lambda)^n\). A fixed-duration collision operation divided by its service time is not automatically this continuous-time model.

\section{New uncertainty transfer when the invariant law changes}
Let \(L_0,L_1\) be symmetric killed conductance Laplacians on the same safe set, and \(M_0,M_1\) positive diagonal invariant mass matrices. Their equilibrium-safe CDFs are
\[
 F_i(t)=\frac{1^TM_i(I-\ee^{-tM_i^{-1}L_i})1}{1^TM_i1}.
\]
If
\[
 \ee^{-\alpha}L_0\preceq L_1\preceq\ee^\alpha L_0,
 \qquad\ee^{-\Delta}M_0\preceq M_1\preceq\ee^\Delta M_0,
\]
then, for every \(t>0\),
\begin{equation}
 \boxed{\ee^{-(\alpha+2\Delta)}F_0(t)
 \le F_1(t)\le\ee^{\alpha+2\Delta}F_0(t).}
 \label{eq:mass}
\end{equation}
Each model uses its own equilibrium-safe initial law. The earlier common-invariant CDF theorem remains the sharper special case.

The new step is a mass monotonicity theorem. Fix \(L\), put \(G=M^{-1}L\), \(u_s=\ee^{-sG}1\), and let
\(H(L,M,t)=1^TM(I-\ee^{-tG})1\). Duhamel differentiation shows that the coefficient of a diagonal mass variation at \(x\) is
\[
 b_x(t)=1-u_t(x)-\int_0^t u_{t-s}(x)(Gu_s)(x)\,ds.
\]
Positivity gives \(0\le u_s\le1\) and \(Gu_s=\ee^{-sG}G1\ge0\), hence \(0\le b_x(t)\le1-u_t(x)\). Unnormalized heat content is coordinatewise increasing in masses. Furthermore
\[
 H(L,cM,t)=cH(L,M,t/c).
\]
Time concavity sandwiches this expression between \(\min(1,c)H\) and \(\max(1,c)H\). The mass band contributes \(\ee^{\pm\Delta}\), the accepted common-form theorem contributes \(\ee^{\pm\alpha}\), and safe normalization contributes another \(\ee^{\pm\Delta}\), proving \eqref{eq:mass}. The discrete version has the explicit, slightly weaker radius
\(B(\alpha)+B(\Delta)+\Delta\), where \(B(z)=2\log((1+\ee^z)/2)\).

\subsection{Acquiring those bands without assuming reversibility under wrong weights}
Under the accepted reciprocal-port model, let local energy errors sum to \(\Delta\), refresh log-error be \(b_0\), and local unknown-support form uncertainty relative to the operating backbone be \(\rho\). The auxiliary conductances must be constructed explicitly. For each unordered local pair, orient \(u\to v\) so fitted \(\pi(u)\le\pi(v)\). Set
\[
 c_{\rm aux}(u,v)=\pi_{\rm fit}(u)r_{\rm true}(u,v),\qquad
 c_{\rm fit}(u,v)=\pi_{\rm fit}(u)\widehat r(u,v).
\]
Define reverse rates by dividing the same symmetric conductance by \(\pi_{\rm fit}(v)\). This ensures reversibility under fitted weights. An absolute forward rate error \(d_\alpha\) gives reverse error no larger than \(d_\alpha\); the accepted backbone bound keeps its constant. True and auxiliary conductances compare by the energy/refresh envelopes. The continuous CDF log-radius becomes
\[
 3\Delta+b_0+\log(1+\rho).
\]
Physically starting in the fitted rather than true equilibrium-safe law adds at most \(2\Delta\). Actual rate/energy intervals need local observations: prepared collision trials, paid at-risk dwell or independently priced energy measurements. Estimating a rare local equilibrium probability by sampling still costs its inverse probability. The result removes inverse \emph{global} rarity from computation, not rare local calibration exposure.

Finite conditional-TV implementation errors have their preserved separate costs. At the exact true equilibrium-safe start, a continuous discrepancy intensity \(v\) yields relative CDF error at most \(v/(\gamma p)\). Arbitrary initial TV error can put mass directly in failure and dominate \(t\downarrow0\). A full physical statement therefore needs exact or multiplicatively controlled safe preparation, or a declared minimum deadline. High-confidence verified factories are substantive supplied interfaces, not an inference from silence about errors.

\section{Reliability evidence, strong alternatives and fixed starts}
The independent diagnostic uses 24 parallel components with local down probability \(1/10\), so stationary failure probability is \(10^{-24}\). A mixture of single and pair refreshes preserves product equilibrium. Permutation symmetry reduces this specific model to 25 states; an ordinary finite-state solver is the better method here and supplies the reference.

The diagnostic actually performed 2,926,345 ordinary kernel draws, using 12,000 trials for each sampled observable. At uniformized deadlines \(0.25,2,20\), the sample-derived intervals contained reference probabilities of about \(1.18\times10^{-25}\), \(9.44\times10^{-25}\), and \(9.44\times10^{-24}\). A reverse window of 448 gave a slow-rate interval
\[
 [4.51077,5.06544]\times10^{-25},
\]
containing the 120-digit numerical reference \(4.7175393\times10^{-25}\). The mean lifetime reference is \(2.1197492\times10^{24}\) uniformized units. The deadline \(10^{24}\) returned \([0.3630581,0.4013994]\), containing the reference \(0.3760930\), without a lifetime-length trajectory.

\begin{figure}[htbp]\centering
\includegraphics[width=\textwidth]{next2_root/figures/reliability_diagnostic.png}
\caption{Executed rare-event diagnostic. The middle points and one tail interval illustrate the operations; they are not an execution of the theorem's generic whole-curve preprocessing. The small symmetric model has a superior ordinary lumped solver.}
\end{figure}

The Palm identity was also checked exactly with rational arithmetic through four updates. Outward arithmetic cases included a \(10^{-100}\) argument and binary deadlines of hundreds of bits. The spectral reference is high-precision numerical diagonalization, not a machine-certified eigenvalue enclosure. Generic Karger counting and the generic whole-curve preprocessing were not executed. Their literal conservative allocation on this model at \(\epsilon=0.1,\delta_{\rm conf}=0.01\) is about \(1.21\times10^{15}\) bounded trials and \(6.21\times10^{19}\) path updates before static, boundary and factory work. These large upper bounds are retained in the evidence.

Karger supplies the static rare-probability machinery \cite{karger}; more recent static algorithms can improve the backend where applicable \cite{cen}. Nakayama--Shahabuddin already combine short-horizon importance sampling and long-horizon regenerative bounds, with broader dynamics and all-working starts \cite{nakayama}. Reversible exponential approximations are established \cite{aldous}. The additional completed object here is a finite nonasymptotic \emph{simultaneous relative certificate} for all deadlines under the stated succinct-model and static-access promises. No historical firstness or general practical superiority is asserted.

A separate root construction completes a chosen-deadline algorithm from the all-working state for \emph{independent} exponential components. A fine time grid preserves first-failure probability multiplicatively: after the first hit, no refresh before the next grid point has probability at least \(\ee^{-A\Delta t}\). Force one failed snapshot drawn from its product marginal, simulate independent component bridges, and divide by the number of failed snapshots. The inverse-multiplicity mean is at least \(1/N_{\rm grid}\). Rational two-state transition approximations preserve every grid path multiplicatively, avoiding an inverse-failure-probability arithmetic cost. At late deadlines, a pointwise product-law burn-in comparison connects the all-working CDF to the equilibrium-safe certificate. The complete note is included; this point-query construction has not been implemented and is not a finitely precomputed all-working curve for all positive times.

Its error scope is necessarily different. Two parallel edges initially working have first-disconnection probability \(\Theta(a^2t^2)\) near zero. An omitted simultaneous failure at rate \(v\) adds \(vt\), causing unbounded relative error as \(t\downarrow0\). The equilibrium-safe intensity theorem cannot be transferred to that start. Neither construction discovers arbitrary hidden common causes or establishes natural-system completeness from short passive data.

\section{Preservation, reproducibility and the remaining boundary}
The archive accompanying this report retains the preceding 47-file deferred-calibration package byte-for-byte, including its nested earlier control-speed and effusion sources. It adds the new mathematical notes, detector construction, executable photon trajectories, explicit resource arithmetic, connected-response scalar diagnostics, complete CDF proof, finite arithmetic code and recorded results. A manifest records relative paths and SHA-256 hashes. Historical incomplete notes remain labelled as such.

The following results remain separate, with no broadened premises:
\begin{itemize}
\item The full classical-feedback all-strength hidden-dephasing acquisition theorem; the sharper global-budget result at its corrected scope; the matched fixed-order parity/GHZ resource law; finite-slew and calibration constructions.
\item The independent thermal-effusion operation, including its thin barrier-shell geometry, actual force/clearing interface, acquisition, timing and remaining exposure cost. It is not repackaged as a cost-free general separation mechanism.
\item Complete output-only Hamiltonian calibration and first-input recovery, the correlated-calibration composition and the distinct bounded-neighborhood no-burn-in construction. The new general passive learner uses paid forgetting; it does not silently inherit the old no-burn-in result.
\item The accepted reliability operator/CDF sensitivity theorems, their sharper shared-backbone constant, fixed-start limitations and finite conditional-update costs. The new whole-CDF algorithm and changing-law result are additional files.
\item P1's completed known-Hamiltonian continuous recovery and bounded individual capture cells, with pure coherent ports, protected memory, growing arrays, precision, propagation and the stated spectral/access-order converses. No fixed aggregate passband or inexpensive microscopic buffer follows.
\end{itemize}

An independent thermodynamic route was also explored. A proposed terminal-information/entropy-production relation does not provide a universal information bound for a freely recorded full trajectory: stationary two-state dynamics already distinguish kinetic rates through jump timing at zero steady entropy production. Its restricted full-machine terminal statement and unfinished construction are retained as incomplete work. It is not counted as a completed consequence.

The present work establishes two complete conditional mathematical capabilities and implements their central observations on synthetic examples. It does not establish an inexpensive general physical reader, an experimentally certified natural-system model or a civilisation-scale deployment. The strongest remaining practical obstruction is now explicit: much tighter attainable inference constants and actual detector/model calibration are needed to make uniform many-body passive learning useful; if arbitrary first-input recovery is required, the coherent archive remains an additional major barrier.

\begin{thebibliography}{99}\small
\bibitem{guta} M. Gu\c{t}\u{a} and J. Kiukas, \emph{Information geometry and local asymptotic normality for multi-parameter estimation of quantum Markov dynamics}, J. Math. Phys. 58, 052201 (2017). \url{https://arxiv.org/abs/1601.04355}.
\bibitem{hagele} D. H\"agele and F. Schefczik, \emph{Higher order moments, cumulants, and spectra of continuous quantum noise measurements}, Phys. Rev. B 98, 205143 (2018). \url{https://arxiv.org/abs/1611.02077}.
\bibitem{sifft} M. Sifft et al., \emph{Quantum polyspectra for modeling and evaluating quantum transport measurements}, Phys. Rev. Research 3, 033123 (2021). \url{https://doi.org/10.1103/PhysRevResearch.3.033123}.
\bibitem{banaszek} K. Banaszek and K. W\'odkiewicz, \emph{Operational Theory of Homodyne Detection}, Phys. Rev. A 55, 3117 (1997). \url{https://arxiv.org/abs/atom-ph/9609001}.
\bibitem{finiteLO} \emph{Statistics of difference events in homodyne detection}, Phys. Rev. A 47, 4227 (1993). \url{https://doi.org/10.1103/PhysRevA.47.4227}.
\bibitem{arad} I. Arad, Z. Chen, N. Guo, P. Rebentrost and Z. Yu, \emph{Near-Optimal Learning of Local Lindbladians}, arXiv:2606.20535v3, revised 3 November 2026. \url{https://arxiv.org/html/2606.20535v3}.
\bibitem{song} J. Song, Y. Zhang, X. Yuan and Y. Wu, \emph{Efficient Lindbladian Learning from Constant-Time Pauli Responses}, arXiv:2607.25795 (2026). \url{https://arxiv.org/abs/2607.25795}.
\bibitem{jimenez} J. Jim\'enez-Rodr\'iguez, G. Franceschetto, A. Ac\'in and L. Pereira, \emph{Rigorous Time-dependent Hamiltonian Learning via Continuous Weak Measurements}, arXiv:2607.16047 (2026). \url{https://arxiv.org/abs/2607.16047}.
\bibitem{karger} D. R. Karger, \emph{A Fully Polynomial Randomized Approximation Scheme for the All Terminal Network Reliability Problem}, SIAM J. Comput. 29, 492--514 (1999); corrected SIAM Review version (2001). \url{https://doi.org/10.1137/S0036144501387141}.
\bibitem{cen} R. Cen, W. He, J. Li and D. Panigrahi, \emph{Beyond the Quadratic Time Barrier for Network Unreliability}, SODA 2024. \url{https://arxiv.org/abs/2304.06552}.
\bibitem{karpluby} R. M. Karp and M. Luby, \emph{A New Monte-Carlo Method for Estimating the Failure Probability of an n-Component System}, UCB/CSD-83-117 (1983). \url{https://www2.eecs.berkeley.edu/Pubs/TechRpts/1983/6351.html}.
\bibitem{aldous} D. Aldous and M. Brown, \emph{Inequalities for rare events in time-reversible Markov chains II}, Stoch. Proc. Appl. 44, 15--25 (1993). \url{https://doi.org/10.1016/0304-4149(93)90035-3}.
\bibitem{nakayama} M. K. Nakayama and P. Shahabuddin, \emph{Quick Simulation Methods for Estimating the Unreliability of Regenerative Models of Large, Highly Reliable Systems}, PEIS 18, 339--368 (2004). \url{https://doi.org/10.1017/S0269964804183058}.
\bibitem{corpus} Accepted Astra scientific corpus, fixed source revision \texttt{2d22b00d4ec6ce4cd894a345c900d53de79ea631}: deferred-calibration and bounded-component-acquisition packages under \texttt{research/returns}. Earlier all-strength, scalable-many-body and unknown-local-port packages remain at revision \texttt{324acb589443002bf8d529cc15e160129c5e36d2}. \url{https://github.com/harjassand/astra-research-knowledge}. Exact package URLs are retained in the accompanying source register.
\end{thebibliography}
\end{document}
